% Results — figure-driven description on the 5 locomotion modes x 3 STDP rates
% ("rehabilitation modes") of the canonical sensory-sited model.
%   fig_stdp_weight_matrix   = weights, both legs x 3 projections x 5 modes
%   fig_stdp_weights_grid    = weights, 5 modes x 3 lambda (3 projections/panel)
%   fig_force_stages_lam*    = force at 3 learning stages, 5 modes, per lambda
%   fig_network_matrix_lam*  = 16 population rates x 5 modes, per lambda
%   fig9_epidural_contrast   = epidural-stim vs natural loading contrast

Before the individual results we fix the terms used throughout. The model is a
two-legged spinal CPG: a spiking network that,
driven only by a weak tonic brainstem input and rhythmic sensory feedback,
produces the alternating extensor/flexor muscle activity of walking. Each leg
contains two mutually-inhibiting neuron pools --- an \emph{extensor} and a
\emph{flexor} rhythm generator (RG-E and RG-F) --- that drive motor neurons and
simple muscle models. A good walking rhythm is one in which extensor and flexor
are active in strict alternation (``counter-phase''); we measure this by the
Pearson correlation between the two muscle forces, which is $-1$ for perfect
alternation and $0$ for none. Three synapses in the model are \emph{plastic} ---
their strengths change as the network runs, through STDP. The Results ask
how this self-tuning builds the gait, and how
well the learned gait survives when sensory feedback is progressively removed,
as it is in spinal-cord-injury rehabilitation.

\subsection{The network tunes its own synapses to a stable set-point}\label{sec:results:weights}

\paragraph{The model and its $5\times5$ grid of test conditions.}
We place the learning in the sensory pathways and hold the brainstem input
fixed, so the three plastic synapses are: the cutaneous CUT$\to$RG-E projection
(paw-contact skin feedback onto the extensor generator) and two proprioceptive
projections, Ia-E$\to$RG-E and Ia-F$\to$RG-F (muscle-stretch feedback onto the
matching generator). Every result is reported on the same grid crossing
\emph{five locomotion modes} with \emph{five learning rates}. The locomotion
modes span a rehabilitation continuum: slow, medium and fast walk (6, 13.5 and
21\,cm/s, all at full weight-bearing --- medium walk is also called
\emph{plantar stepping} or \emph{baseline}), then \emph{toe stepping} (partial
body-weight support) and \emph{air stepping} (the limb fully unloaded, as when a
spinal rat is suspended). The learning rate $\lambda$ sets how much each plastic
synapse changes per correlated pre/post spike pair --- larger $\lambda$ means
faster learning --- and we sweep it across five decades,
$\lambda\in\{10^{-2},10^{-3},10^{-4},10^{-5},10^{-6}\}$.

\paragraph{The synapses grow, then settle at values that do not depend on the
walking mode.}
Starting from weak random values, STDP drives each plastic synapse to a steady
strength and holds it there (Figure~\ref{fig:weight_matrix}). The cutaneous
CUT$\to$RG-E synapse grows strong --- to $\approx63$\,pA --- in every
weight-bearing mode, whereas the two proprioceptive synapses settle at a low
$\approx4$--$5$\,pA, well below their $10$\,pA ceiling. This low set-point is
functionally important: a light, correctly-timed proprioceptive input reinforces
each burst without ``filling in'' the silent gap between bursts, which would
otherwise abolish the alternation. The left and right legs learn identical
values.

\begin{figure}[t]
  \centering
  \includegraphics[width=0.92\textwidth]{fig_stdp_weight_matrix.png}
  \caption{\textbf{STDP weight trajectories of the three plastic projections
    across the five locomotion modes.} Rows: CUT$\to$RG-E, Ia-E$\to$RG-E,
    Ia-F$\to$RG-F for the left leg (white) then the right leg (tinted); columns:
    slow~/ medium\,$=$\,plantar\,$=$\,baseline / fast walk and toe / air
    stepping. Each panel overlays the three learning rates $\lambda$. Weight
    (pA) versus time (s). CUT$\to$RG-E converges to $\approx63$\,pA in every
    mode; the two proprioceptive projections self-stabilise at a low
    $\approx4$--$5$\,pA set-point; the two legs are symmetric.}
  \label{fig:weight_matrix}
\end{figure}

\paragraph{The learning rate sets only how fast the synapse settles; unloading
slows it.}
Faster learning reaches the plateau sooner (Figure~\ref{fig:weights_grid}):
CUT$\to$RG-E settles within $\approx7$\,s at $\lambda=10^{-3}$, but takes
$\approx75$\,s at $10^{-4}$ and has not finished settling within the $120$\,s run
at $10^{-5}$. Removing body-weight support slows learning further, because the
paw-contact feedback that drives the synapse is itself weakened when the limb is
unloaded: under air stepping even the fastest rate climbs only gradually (to
$\approx50$\,pA by the end of the run). The proprioceptive synapses stay at
their low set-point in every condition.

\begin{figure}[t]
  \centering
  \includegraphics[width=0.92\textwidth]{fig_stdp_weights_grid.png}
  \caption{\textbf{Weight trajectories of all three plastic projections, per
    rehabilitation mode.} Rows: the five locomotion modes; columns: the three
    learning rates. Each panel shows CUT$\to$RG-E (solid, left axis, pA) and the
    two proprioceptive projections Ia-E$\to$RG-E and Ia-F$\to$RG-F (dashed, right
    axis, pA) versus time (s). Panels (a)--(y). The learning rate sets the
    time-to-plateau; unloading (bottom rows) slows CUT$\to$RG-E convergence.}
  \label{fig:weights_grid}
\end{figure}

% =============================================================
\subsection{Rehabilitation progress: force develops as the weight converges}\label{sec:results:force}

We next examine whether this self-tuned wiring produces good walking, and how
the gait builds up over a session. Figures~\ref{fig:force_l2}--\ref{fig:force_l6} show the
extensor and flexor muscle forces in three short windows --- the beginning,
middle and end of each $120$\,s run --- so that reading a figure left-to-right
traces how the gait develops as the synapse strengthens (one figure per learning
rate). Early in the run the two forces overlap in ragged, poorly-timed bursts;
as CUT$\to$RG-E grows they separate into a clean alternation. At the fast and
intermediate rates ($\lambda=10^{-2}$ to $10^{-4}$) the synapse converges within
the run and the weight-bearing modes reach near-perfect counter-phase (muscle-%
force correlation down to $-0.99$, where $-1$ is ideal). Toe stepping is clean at
the faster rates but degrades as learning slows (correlation $-0.96$ at $10^{-2}$
down to $-0.39$ at $10^{-6}$), and air stepping never exceeds $\approx-0.55$ at
any rate, because with the limb unloaded the extensor synapse stays weak
($\approx4$\,pA). At the two slowest rates ($\lambda=10^{-5},10^{-6}$) the synapse
barely develops within $120$\,s, so the gait stays under-formed from beginning to
end.

\begin{figure}[t]
  \centering
  \includegraphics[width=0.9\textwidth]{fig_force_stages_lam1em2.png}
  \caption{\textbf{Force at three learning stages, five locomotion modes,
    $\lambda=10^{-2}$ (fastest rate).} Rows: the five locomotion modes; columns:
    beginning / middle / end of the $120$\,s run. Solid = Force-E, dashed =
    Force-F (a.u.)\ versus time (s); panels annotated with the mean CUT$\to$RG-E
    weight and in-window corr$(F_E,F_F)$. Panels (a)--(o).}
  \label{fig:force_l2}
\end{figure}

\begin{figure}[t]
  \centering
  \includegraphics[width=0.9\textwidth]{fig_force_stages_lam1em3.png}
  \caption{\textbf{Force at three learning stages, five locomotion modes,
    $\lambda=10^{-3}$.} Same layout as Figure~\ref{fig:force_l2}.}
  \label{fig:force_l3}
\end{figure}

\begin{figure}[t]
  \centering
  \includegraphics[width=0.9\textwidth]{fig_force_stages_lam1em4.png}
  \caption{\textbf{Force at three learning stages, five locomotion modes,
    $\lambda=10^{-4}$.} Same layout as Figure~\ref{fig:force_l3}. At this
    intermediate rate the three stages are most distinct, spanning the full
    $120$\,s convergence.}
  \label{fig:force_l4}
\end{figure}

\begin{figure}[t]
  \centering
  \includegraphics[width=0.9\textwidth]{fig_force_stages_lam1em5.png}
  \caption{\textbf{Force at three learning stages, five locomotion modes,
    $\lambda=10^{-5}$.} Same layout as Figure~\ref{fig:force_l2}. At this slow
    rate the descending weight barely develops within the run, so the gait stays
    under-developed across all three windows.}
  \label{fig:force_l5}
\end{figure}

\begin{figure}[t]
  \centering
  \includegraphics[width=0.9\textwidth]{fig_force_stages_lam1em6.png}
  \caption{\textbf{Force at three learning stages, five locomotion modes,
    $\lambda=10^{-6}$ (slowest rate).} Same layout as Figure~\ref{fig:force_l2};
    the descending weight remains near its initial value throughout, so
    counter-phase is carried almost entirely by the paced core.}
  \label{fig:force_l6}
\end{figure}

% =============================================================
\subsection{Network-level activity across the locomotion modes}\label{sec:results:network}

The muscle force is only the read-out of a whole network, so we next ask whether
the clean alternation is a property of the entire circuit.
Figures~\ref{fig:net_l2}--\ref{fig:net_l6} show the firing rates of all sixteen
recorded neuron pools --- for both legs, the rhythm generators (RG-E, RG-F), the
reciprocal-inhibition interneurons (InE, InF), the proprioceptive interneurons
(IaInt-E, IaInt-F) and the motor pools (M-E, M-F) --- one figure per learning
rate. In slow, medium and fast walk and in toe stepping, every extensor-side
pool fires in strict alternation with its flexor-side partner, the two legs stay
half a cycle apart (a $180^{\circ}$ offset, as in a real trot), and the
inhibitory interneurons alternate just as cleanly as the rhythm generators --- so
the alternation is built into the core circuit, not imposed only at the muscle.
The exception is air stepping (rightmost column): with the sensory feedback
withdrawn, the extensor-side pools (RG-E, InE, IaInt-E, M-E) fall nearly silent
while the flexor side, which bursts on its own intrinsic dynamics, keeps going
--- the network-level signature of the walking failure seen in the muscle forces.

\begin{figure}[t]
  \centering
  \includegraphics[width=0.95\textwidth]{fig_network_matrix_lam1em2.png}
  \caption{\textbf{Full-circuit population activity across the five locomotion
    modes, $\lambda=10^{-2}$ (fastest rate).} Rows: the sixteen recorded
    populations (both legs; right-leg rows tinted); columns (a)--(e): slow /
    medium=plantar=baseline / fast walk and toe / air stepping. Population rate
    (spikes$\cdot$neuron$^{-1}\cdot$s$^{-1}$) versus time (ms), last cycles.}
  \label{fig:net_l2}
\end{figure}

\begin{figure}[t]
  \centering
  \includegraphics[width=0.95\textwidth]{fig_network_matrix_lam1em3.png}
  \caption{\textbf{Full-circuit population activity, $\lambda=10^{-3}$.} Same
    layout as Figure~\ref{fig:net_l2}.}
  \label{fig:net_l3}
\end{figure}

\begin{figure}[t]
  \centering
  \includegraphics[width=0.95\textwidth]{fig_network_matrix_lam1em4.png}
  \caption{\textbf{Full-circuit population activity, $\lambda=10^{-4}$.} Same
    layout as Figure~\ref{fig:net_l3}.}
  \label{fig:net_l4}
\end{figure}

\begin{figure}[t]
  \centering
  \includegraphics[width=0.95\textwidth]{fig_network_matrix_lam1em5.png}
  \caption{\textbf{Full-circuit population activity, $\lambda=10^{-5}$.} Same
    layout as Figure~\ref{fig:net_l2}. At this under-converged rate the
    counter-phase is weaker, most visibly at the higher cadences.}
  \label{fig:net_l5}
\end{figure}

\begin{figure}[t]
  \centering
  \includegraphics[width=0.95\textwidth]{fig_network_matrix_lam1em6.png}
  \caption{\textbf{Full-circuit population activity, $\lambda=10^{-6}$
    (slowest rate).} Same layout as Figure~\ref{fig:net_l2}.}
  \label{fig:net_l6}
\end{figure}

% =============================================================
\subsection{Comparison across the 5 locomotion modes $\times$ 5 STDP rates}\label{sec:results:comparison}

To compare all twenty-five conditions at a glance, Figures~\ref{fig:mode_lambda_heat}
and~\ref{fig:mode_lambda_trend} and Table~\ref{tab:mode_lambda} reduce each
condition to four numbers (measured over the last $20$\,s of the run): the
alternation of the two muscle forces [corr$(F_E,F_F)$] and of the two rhythm
generators [corr(RG-E,RG-F)] --- both near $-1$ when alternation is clean --- the
peak extensor force, and the final strength of the learned CUT$\to$RG-E synapse.
Three patterns stand out. \textbf{(i) A ``sweet spot'' for the learning rate.}
The final synaptic strength grows steadily with $\lambda$ --- from the
$\approx4$\,pA floor at $10^{-6}$ up to $\approx64$\,pA for $10^{-4}$ and faster
--- but the \emph{quality} of the alternation does not follow suit: it is best at
the intermediate rates and slightly noisier at the very fastest rate, giving a
broad optimum around $\lambda=10^{-3}$--$10^{-4}$. \textbf{(ii) A loading
gradient.} At every learning rate the alternation and the peak force fall as the
limb is unloaded, and under air stepping the muscle-force alternation settles
near $-0.5$ regardless of $\lambda$ --- the withdrawn feedback caps how much
descending drive the network can build. \textbf{(iii) Robustness to walking
speed.} Among the weight-bearing modes the four numbers are almost flat across
speed once learning has converged, showing that the sensory-paced core carries
the gait across the whole rat locomotor range.

\paragraph{Counter-phase \emph{recovers} from toe to air stepping at the
slow learning rates.}
A counterintuitive feature of Figure~\ref{fig:mode_lambda_trend} is that at
the under-converged rates ($\lambda=10^{-5},10^{-6}$) both correlations do not
decrease monotonically with unloading: they are \emph{weakest at toe stepping}
and then \emph{strengthen} again at air stepping (e.g.\ at $\lambda=10^{-6}$,
corr$(F_E,F_F)$ goes $-0.39$ at toe $\to -0.51$ at air, and corr(RG) $-0.74\to
-0.81$), whereas at the fast rates the correlations fall monotonically. This is
the same non-monotonic ``intermediate drive is worst'' effect seen in the
descending-arm learning-rate dependence and in the frozen-weight control: a
\emph{partially}-developed CUT$\to$RG-E weight (toe stepping, where the gated
afferents leave the projection at an intermediate $\approx7$--$18$\,pA) is
strong enough to perturb the paced reciprocal-inhibition core but too weak and
too poorly-timed to cleanly reinforce it, so it desynchronises the rhythm. Under
air stepping the afferents are withdrawn almost entirely (CUT$\to$RG-E gated to
$\approx4$--$8$\,pA), removing that interfering intermediate drive, and the
rhythm reverts to the clean counter-phase \emph{floor} imposed by the paced
heel$\to$toe drive and reciprocal inhibition alone. The paced core therefore
sets a lower bound on counter-phase that the mid-loading, mid-learning regime
can actually fall below.

\paragraph{The RG correlation stays stronger than the force correlation.}
Across the whole matrix corr(RG-E,RG-F) is consistently more negative than
corr$(F_E,F_F)$ --- most visibly under air stepping ($\approx-0.8$ versus
$\approx-0.5$). The rhythm generators keep alternating there because their
antiphase is produced by reciprocal inhibition plus the external pacing, neither
of which requires the plastic drive; what fails under unloading is the
\emph{muscle-force expression} of the extensor. Because the only large plastic
projection is extensor-biased (CUT$\to$RG-E, with no flexor counterpart), gating
the afferents starves the extensor of the activation it needs to develop force,
so the extensor force collapses while the flexor --- intrinsically bursting ---
persists. The counter-phase is thus lost at the muscle level before it is lost
at the rhythm-generator level.

\begin{figure}[t]
  \centering
  \includegraphics[width=0.95\textwidth]{fig_mode_lambda_heatmaps.png}
  \caption{\textbf{Metric heatmaps across the $5\times5$ rehabilitation-mode
    matrix} (rows: locomotion modes; columns: STDP rate $\lambda$). corr$(F_E,F_F)$
    and corr(RG-E,RG-F) (diverging scale), peak Force-E (a.u.), and converged
    CUT$\to$RG-E weight (pA). Last-20\,s window, canonical sensory model.}
  \label{fig:mode_lambda_heat}
\end{figure}

\begin{figure}[t]
  \centering
  \includegraphics[width=0.9\textwidth]{fig_mode_lambda_trends.png}
  \caption{\textbf{Metric trends across the locomotion modes, one line per STDP
    rate $\lambda$.} Same four metrics as Figure~\ref{fig:mode_lambda_heat},
    plotted against the locomotion mode (slow walk $\to$ air stepping) to expose
    the loading gradient and the learning-rate ordering.}
  \label{fig:mode_lambda_trend}
\end{figure}

\input{mode_lambda_table}

% =============================================================
\subsection{Epidural stimulation rescues stepping under unloading}\label{sec:results:epidural}

Finally, the model speaks to a striking clinical observation: people and animals
with severe spinal cord injury often cannot step on their own, yet regain
stepping when the lumbar spinal cord is electrically stimulated (\emph{epidural
stimulation}). We reproduce the two situations by choosing how the paw-contact
drive behaves under unloading (Figure~\ref{fig:epidural}). In the \emph{natural}
(unassisted) case both the cutaneous and the proprioceptive feedback fade
together as the limb is unloaded, and the rhythm collapses (muscle-force
correlation $-0.41$ at air stepping, with the cycle period also becoming
irregular). When the rhythmic cutaneous drive is instead held on regardless of
loading --- the model's stand-in for epidural stimulation, which supplies a
rhythmic input that does not depend on limb
contact~\cite{courtine2009,harkema2011} --- the alternation is preserved across
the whole loading range (correlation $-0.94$ at baseline to $-0.97$ at air
stepping). This double dissociation
captures the mechanism behind the clinical finding that motor-complete subjects
cannot step unassisted but recover treadmill stepping under epidural
stimulation~\cite{courtine2009,harkema2011}.

\begin{figure}[t]
  \centering
  \includegraphics[width=0.8\textwidth]{fig9_epidural_contrast.png}
  \caption{\textbf{Epidural-stimulation rescue --- paced cutaneous drive
    preserves stepping under unloading.} Loading swept from baseline to air
    stepping. \textbf{(a)} counter-phase corr$(F_E,F_F)$ versus loading:
    epidural-stim (CUT held intact) stays strong, the natural (CUT gated) arm
    collapses. \textbf{(b)} cycle-period jitter versus loading.
    \textbf{(c, d)} representative air-stepping force traces for the stim and
    natural arms.}
  \label{fig:epidural}
\end{figure}

% =============================================================
\subsection{Closed-loop force-triggered stepping and tag-and-capture
  consolidation}\label{sec:results:consolidate}

\emph{Debug-scale, preliminary: a full production-scale (N=100, brainstem
input at 60\,Hz) confirmatory sweep has been prepared for submission to an
HPC cluster, and every figure in this subsection will be replaced with the
production-scale equivalent once it completes; the qualitative pattern
reported here is not expected to change, but the numbers are debug-scale
(reduced population size, 20\,Hz brainstem input) and should be read
accordingly.}

The results above drive stance/swing timing with an explicit clock. We also
built a closed-loop alternative in which cutaneous (paw-contact) firing is
gated directly on each leg's own extensor force: contact turns on when force
rises through 80\% of that stride's peak and turns off when it falls back
through 35\%, so the rhythm is generated by the mechanics of stepping itself
rather than imposed on it. Because the extensor rhythm generator has no
intrinsic burst-terminating current, a fixed-duration failsafe (matched to
the model's own $\approx450$\,ms stance target) forces the phase transition
if the force-threshold crossing is late; we track the fraction of strides
that hit this failsafe (`frac\_at\_cap`) as a direct check on whether a given
configuration is a genuine closed loop or a disguised clock wearing a
force-shaped costume.

\paragraph{Five locomotion modes, two learning architectures.} We confirmed
two additional timing points bracketing the original one --- a
$\approx31\%$ slower and a $\approx13$--$18\%$ faster variant --- found by
scaling the stance-duration and muscle-fatigue time constants together
rather than moving any one parameter alone. The faster point needed one
further ingredient: a small, \emph{persistent} left/right asymmetry in the
muscle-fatigue time constant (the leading leg fatigues slightly faster
than the other, for the whole run, not just at the initial symmetry-breaking
step). Without it, several otherwise-genuine fast configurations settled
into a perfectly regular cycle whose left/right phase relationship
(alternating vs.\ moving together) was decided by run-to-run noise rather
than by any timing parameter; the persistent asymmetry --- itself
bio-plausible, since real limbs are not identical --- reliably resolves
this bistability. We crossed both walking-speed points, plus the original
medium speed, with the two ablation loading levels used throughout this
paper (toe and air stepping, partial and full unloading) and the two
learning architectures: a \emph{descending} arm in which the
brainstem-to-rhythm-generator synapse remains plastic (as in the results
above), and a \emph{sensory} arm in which that synapse is frozen at a weak
fixed value and only the cutaneous and proprioceptive afferent synapses
learn --- the model's stand-in for a rehabilitation scenario in which
descending drive is compromised but peripheral feedback pathways are intact.

Figure~\ref{fig:ct_all_modes_force} shows all three walking speeds and toe
stepping produce a genuine, closed-loop counter-phase rhythm in both
architectures (corr$(F_E,F_F)=-0.53$ to $-0.90$ at the three speeds,
$-0.23$ to $-0.84$ at toe stepping, weakening but not collapsing under
partial unloading as expected). Air stepping is the exception, and
deliberately so: force stays too small for the closed-loop threshold
mechanism to find any genuine crossings, so cutaneous drive never
develops (Figure~\ref{fig:ct_all_modes_weights}, right column) and the
rhythm degrades toward the tick resolution of the simulation --- the
same natural-air-stepping collapse already reported in
Section~\ref{sec:results:epidural} under the fixed clock, now reproduced
under closed-loop timing instead of imposed by construction. We do not
attempt to force a clean rhythm out of natural air stepping here; the
model's own epidural-stimulation rescue (Section~\ref{sec:results:epidural})
is the appropriate comparison condition for that, not a retuned threshold.
Figure~\ref{fig:ct_all_modes_weights} shows the CUT$\to$RG-E weight
converging to the same $\approx60$--$67$\,pA plateau at every walking speed
and at toe stepping, in both architectures --- the same
initialization-independent set-point reported in
Section~\ref{sec:results:weights}, now shown to hold under closed-loop
timing as well as the fixed clock, across a substantially wider range of
speeds and loading than the single operating point first confirmed for
this mechanism.

\begin{figure}[t]
  \centering
  \includegraphics[width=0.92\textwidth]{fig_forcetrigger_all_modes_force_stages.png}
  \caption{\textbf{Closed-loop (force-triggered) stepping across all five
    locomotion modes and both learning architectures.} Rows (top to
    bottom): descending arm at slow/medium/fast/toe/air, then sensory arm
    at slow/medium/fast/toe/air; columns: beginning / middle / end of the
    $60$\,s run, as in Figure~\ref{fig:force_l2}. Air stepping (rows 5, 10)
    shows the expected natural collapse, not a fitting failure. Panels
    (a)--(ad).}
  \label{fig:ct_all_modes_force}
\end{figure}

\begin{figure}[t]
  \centering
  \includegraphics[width=0.98\textwidth]{fig_forcetrigger_all_modes_weights_grid.png}
  \caption{\textbf{The learned set-point is the same across speeds and
    loading levels, except where force never develops.} CUT$\to$RG-E
    weight (solid, left axis; dotted = captured baseline where
    consolidation, below, is active --- medium, fast and toe rows, all at
    the same confirmed gain; slow and air run without consolidation, see
    text) and Ia-E/F$\to$RG weights (dashed, right axis) versus time, one
    panel per condition from Figure~\ref{fig:ct_all_modes_force}. Toe
    stepping's Ia$\to$RG weights climb further than at full loading (a
    loading-dependent relaxation of the Ia$\to$RG weight ceiling lets
    proprioceptive feedback take over more excitatory drive as cutaneous
    input is progressively removed); air stepping's CUT$\to$RG-E rises only
    slowly and by a small amount over the run, well short of the
    $\approx60$--$67$\,pA plateau reached elsewhere.}
  \label{fig:ct_all_modes_weights}
\end{figure}

\paragraph{Tag-and-capture consolidation.} Vanilla STDP retains every
potentiation permanently, up to its ceiling, with no way to distinguish a
genuine, force-threshold-driven stride from one only completed because the
failsafe forced it. We therefore added an optional consolidation layer,
directly analogous to synaptic tagging-and-capture in hippocampal memory
models \todo{cite Sandkuhler spinal E-LTP/L-LTP; Grau contingency-gated
spinal learning; Wolpaw two-phase H-reflex conditioning}: each plastic
connection keeps its native STDP weight (the fast, local ``tag'') alongside
a separate, slower-changing \emph{baseline}. A shared per-leg pool
accumulates from genuine stride endings and is drawn down by failsafe-forced
ones; only once it crosses a threshold does the baseline snap to the current
weight (a \emph{capture} event), and between captures the live weight
continuously relaxes back toward the baseline. The weight ceiling itself is
untouched --- this mechanism governs what is \emph{retained} within the
existing ceiling, not the ceiling.

\todo{Revision note, to resolve before submission: the medium-speed
timing point used to illustrate this mechanism was found to have 4 of its
5 timing constants landing exactly on the simulation's discretization
step, which let its stance/swing timing collapse onto seed- and
thread-sensitive discrete states independently of consolidation. A small
($\approx1\%$) de-alignment of two constants fixed this for the base
(non-consolidated) mechanism. Separately, the bookkeeping for the
brainstem-to-rhythm-generator pathway was tracked but never written back
to the simulator, so that pathway ran unconsolidated (ordinary STDP)
regardless of this mechanism being enabled, for every result reported
below at the time this note was written. Both issues are now fixed, and
the gain search below was re-run afterward; the numbers and figures in
this subsection reflect the corrected code, but only at reduced network
size (debug scale) --- a full-scale confirmatory run is in progress.
Remove this note once that confirmation lands and the numbers below are
verified to match at full scale.}

One gain setting (how strongly a genuine vs.\ a forced stride moves the
pool) was identified by grid search and confirmed across two random
seeds, and applies uniformly across all three walking speeds and both
learning architectures: $0.25$ for a genuine stride, $0.10$ for a
failsafe-forced one. This supersedes an earlier finding of two
architecture-specific settings, an artifact of the two issues noted above
rather than a genuine property of the two learning architectures.
Figure~\ref{fig:consolidate_desc} shows the descending arm:
without consolidation the weight is retained unconditionally regardless of
how it was earned; at a forced-favoring gain (the reverse ratio, forced
outweighing genuine 3:1) the pool never crosses threshold and the
mechanism is an inert no-op --- notably, the brainstem-to-rhythm-generator
pathway is visibly suppressed here too (converging far below its
unconsolidated plateau) even though it never captures, since the
continuous tag-decay this mechanism applies runs whether or not a capture
ever happens; at the confirmed gain the baseline visibly steps upward each
time enough genuine strides accumulate
(Figure~\ref{fig:consolidate_desc_weights}), and all three plastic
pathways --- including the brainstem projection --- converge to a stable
captured value.

\begin{figure}[t]
  \centering
  \includegraphics[width=0.92\textwidth]{fig_consolidate_descending_force_stages.png}
  \caption{\textbf{Tag-and-capture consolidation, descending arm.} Rows:
    no consolidation; an inert (never-captures) gain setting ($0.15$
    genuine / $0.30$ forced); the confirmed working gain ($0.25$ genuine /
    $0.10$ forced). Columns as in Figure~\ref{fig:ct_all_modes_force}.
    Panels (a)--(i).}
  \label{fig:consolidate_desc}
\end{figure}

\begin{figure}[t]
  \centering
  \includegraphics[width=0.95\textwidth]{fig_consolidate_descending_weights_grid.png}
  \caption{\textbf{The capture staircase, now spanning all three plastic
    pathways.} Solid/dashed = live STDP weight; dotted = captured baseline
    (only present once consolidation is on). Left: no consolidation (no
    baseline to show). Middle: inert setting --- baseline frozen at
    initialization despite the live weight moving, on every pathway
    including the brainstem projection. Right: confirmed working gain ---
    baseline steps upward at each capture event and tracks the live
    weight's converged value by the end of the run, again on all three
    pathways.}
  \label{fig:consolidate_desc_weights}
\end{figure}

\paragraph{The same gain also works for the sensory arm.} Earlier testing,
before the two fixes noted above, had found the descending arm's gain
setting synchronized the legs when applied to the sensory arm instead, and
concluded the two architectures needed different settings. That
conclusion does not hold once the two fixes are applied: the confirmed
gain (Figure~\ref{fig:consolidate_sens}) captures cleanly in the sensory
arm too, converging both legs' cutaneous synapse to a captured baseline
(Figure~\ref{fig:consolidate_sens_weights}) without reproducing the
earlier synchronization failure. \todo{Revision note: a bout-timeline figure previously illustrated a
settling-in transient at run start (an extended run of failsafe-forced
strides early on that resolves to genuine timing before the run ends,
motivating the steady-state-only convention used throughout this
subsection). That figure and claim are removed here pending
re-verification: the operating point used to demonstrate it has changed
along with everything else in this subsection, and a first check on the
new fast-timing point found no such transient (steady from the start), so
the original claim should not be carried forward uncritically. Re-add once
a current, verified version exists.}

\begin{figure}[t]
  \centering
  \includegraphics[width=0.92\textwidth]{fig_consolidate_sensory_force_stages.png}
  \caption{\textbf{Tag-and-capture consolidation, sensory arm.} Rows: no
    consolidation; the confirmed working gain ($0.25$ genuine / $0.10$
    forced --- the same setting as the descending arm). Columns as in
    Figure~\ref{fig:ct_all_modes_force}. Panels (a)--(f).}
  \label{fig:consolidate_sens}
\end{figure}

\begin{figure}[t]
  \centering
  \includegraphics[width=0.95\textwidth]{fig_consolidate_sensory_weights_grid.png}
  \caption{\textbf{Same layout as Figure~\ref{fig:consolidate_desc_weights},
    sensory arm.} Left: no consolidation. Right: confirmed gain --- the
    cutaneous synapse's baseline steps upward and tracks the converged
    live weight, the same capture behavior seen in the descending arm.}
  \label{fig:consolidate_sens_weights}
\end{figure}

\paragraph{The same gain extends to the faster and toe-stepping timing
points once one further time constant is rescaled.} The gain search above
was run at the medium-speed timing point; naively applying the same
setting at the faster point or at toe stepping reproduced failure
signatures resembling those first documented for the sensory arm above
(failsafe domination or synchronization) even where the underlying
closed-loop timing itself was already genuine. The cause was a fixed leak
time constant governing how quickly an uncaptured tag relaxes back toward
its baseline: consolidation applies this leak continuously, not only
between captures, so a value tuned for the medium-speed cycle length
actively suppressed weight growth at the shorter cycle lengths of the
faster and toe-stepping points, independently of the gain. Rescaling this
one constant (roughly in proportion to how much faster the local cycle
is) restores the confirmed gain to genuine operation at both new points,
in both architectures. The three walking speeds are not uniformly clean,
however: fast reproduces as a near-identical cross-seed match (steady-state
corr$(F_{E,L},F_{E,R})$ within $0.001$ of itself across seeds), while toe's
anti-phase signal, though same-signed across seeds, is markedly weaker and
leaves a small residual failsafe-domination fraction on one leg --- a
genuine but less robust confirmation than fast's. The converged weight
trajectories for all three walking speeds are the ``medium'', ``fast'' and
``toe'' columns of Figure~\ref{fig:ct_all_modes_weights} above. The slow
walking point is the one exception found so far: the same rescaling
narrows but does not eliminate a residual left/right synchronization
tendency under consolidation, so Figure~\ref{fig:ct_all_modes_force}'s slow
rows run without consolidation (STDP alone, which is already excellent
there on its own) rather than with a setting confirmed to occasionally
misbehave.

\paragraph{Scope of this preliminary result.} This subsection now covers
four of the five locomotion modes under closed-loop timing (slow, medium,
fast walking and toe stepping), both learning architectures, and
tag-and-capture consolidation confirmed at three of those four (medium,
fast, toe) with a single shared gain. Slow alone runs consolidation-free,
by choice: the base mechanism is already excellent there on its own, and
consolidation was found actively harmful across a wide gain/threshold
search. Air stepping is deliberately excluded from the closed-loop-timing
and consolidation results, not because the mechanism fails to fit it, but
because a genuinely unloaded limb should not produce a clean closed-loop
rhythm at all --- see the discussion above and
Section~\ref{sec:results:epidural}. \todo{Revision note: everything in
this subsection, including the confirmed gain, is currently validated only
at reduced network size (debug scale); a full-scale confirmatory sweep
covering medium, fast and toe, both architectures, with consolidation
enabled at the confirmed gain, is in progress. Update this paragraph and
the preceding one once it lands --- do not cite the numbers above as
production-scale confirmed until then.}
